
\chapter{Introdução}
\label{cap:introducao}

A Engenharia de Requisitos (ER) é uma fase fundamental do desenvolvimento de software, responsável por estabelecer a ponte entre as necessidades das partes interessadas e a implementação técnica do sistema. Na maioria das organizações, os requisitos são levantados em reuniões conduzidas por profissionais especializados (analistas de requisitos, \textit{product owners} ou \textit{scrum masters}), cujo papel é capturar com precisão discussões complexas e em constante evolução e transformá-las em artefatos formais \cite{bokhari2011metrics}. Esse fluxo de trabalho manual introduz riscos amplamente conhecidos: a perda de informação causada pelo intervalo de tempo entre a reunião e a etapa de documentação, a sobrecarga cognitiva do analista, que precisa simultaneamente participar da discussão e registrá-la, e a inconsistência de granularidade e de formato entre os artefatos resultantes. Como os documentos de requisitos servem de base para todas as decisões subsequentes do desenvolvimento, especificações incompletas ou de baixa qualidade traduzem-se diretamente em implementações desalinhadas, retrabalho e estouro de custos.

Ferramentas comerciais de inteligência de reuniões (como Fellow.ai \cite{fellow2025}, Jamie \cite{jamie2025} e Otter.ai) automatizam a transcrição e a extração de itens de ação, mas produzem anotações de propósito geral, e não artefatos formais de requisitos de software. A tradução da discussão de uma reunião para uma Especificação de Requisitos de Software (do inglês \textit{Software Requirements Specification}, SRS) estruturada permanece uma etapa manual e intensiva em conhecimento especializado. Na literatura acadêmica, os estudos que aplicam Modelos de Linguagem de Grande Escala (do inglês \textit{Large Language Models}, LLMs) à ER tratam predominantemente de subtarefas isoladas: a geração de histórias de usuário a partir de requisitos textuais já existentes \cite{delgado2024aidriven}, a extração de modelos conceituais \cite{ren2024combining} ou a classificação de requisitos. Poucos trabalhos exploram \textit{pipelines} completos, de ponta a ponta, que partam do áudio bruto de uma reunião e cheguem a uma documentação formalmente estruturada e específica para as partes interessadas \cite{motger2025share}.

Este trabalho aborda essa lacuna por meio de uma \emph{arquitetura multiagente} na qual quatro agentes especializados colaboram para transformar o áudio de uma reunião em documentação de requisitos. Decompor o \textit{pipeline} em agentes, em vez de utilizar um único \textit{prompt} monolítico de LLM, oferece vantagens concretas de engenharia: (i)~\textit{separação de responsabilidades}, permitindo que cada agente seja desenvolvido, testado e evoluído de forma independente; (ii)~\textit{flexibilidade de ponto de entrada}, permitindo que o usuário ingresse no \textit{pipeline} em qualquer estágio (por exemplo, a partir de uma gravação ou de uma transcrição já existente); (iii)~\textit{seleção de modelo por agente}, ajustando a capacidade e o custo do LLM à complexidade de cada tarefa; e (iv)~\textit{propagação de erro observável}, tornando possível rastrear e medir como os erros de transcrição se propagam para os artefatos a jusante, uma preocupação de engenharia de primeira ordem em sistemas agênticos.

Experimentos iniciais com os Agentes~1 e~2, conduzidos em uma etapa anterior desta pesquisa, demonstraram 100\% de acurácia na identificação de personas (28 de 28 personas ao longo de 7 reuniões do \textit{corpus} AMI) e 83\% de conformidade com o critério INVEST nas histórias de usuário geradas. Esta monografia estende aquela implementação à sua \emph{forma completa}, acrescentando os Agentes~3 e~4 para a geração do documento de requisitos e a síntese do diagrama de domínio, introduzindo dois formatos de saída selecionáveis (a SRS no padrão IEEE~830 e um formato específico de cliente) e conduzindo uma validação no mundo real sobre um projeto de software real de um órgão público. Substitui-se, ainda, a avaliação qualitativa preliminar por um conjunto de métricas quantitativas que cobre a acurácia de transcrição, métricas de extração de informação, uma pesquisa de opinião com profissionais e a análise da propagação de erro ao longo do \textit{pipeline}.

\section{Objetivos}

O objetivo geral deste trabalho é desenvolver e avaliar um sistema multiagente, baseado em inteligência artificial generativa, capaz de automatizar a elicitação de requisitos de software a partir do áudio de reuniões, produzindo documentos de requisitos estruturados.

Para atingir esse objetivo geral, foram definidos os seguintes objetivos específicos:

\begin{enumerate}
    \item[1.] Implementar um \textit{pipeline} de transcrição automática de áudio de reuniões em português, robusto a arquivos longos e a diferentes formatos de mídia.
    \item[2.] Identificar automaticamente as personas (atores do sistema) presentes na reunião e gerar histórias de usuário segundo o critério INVEST.
    \item[3.] Gerar documentos de requisitos estruturados em dois formatos distintos (o padrão IEEE~830 e um formato específico de cliente), a partir das informações extraídas.
    \item[4.] Validar o sistema em um caso real, por meio de uma parceria com a equipe de tecnologia de um órgão público, comparando o documento gerado com um documento oficial produzido manualmente (\textit{gold standard}).
    \item[5.] Avaliar quantitativamente o desempenho de cada agente e analisar empiricamente a propagação de erros entre os estágios do \textit{pipeline}.
\end{enumerate}

\section{Questões de Pesquisa}

A partir dos objetivos delineados, este trabalho busca responder às seguintes questões de pesquisa (QP):

\begin{enumerate}
    \item[\textbf{QP1.}] Qual é a acurácia da transcrição automática (WER/CER) do Agente~1 para áudio de reuniões em português?
    \item[\textbf{QP2.}] Com que precisão o Agente~2 identifica as personas das partes interessadas e gera histórias de usuário (Precisão/Revocação/F1 e conformidade com o INVEST)?
    \item[\textbf{QP3.}] Qual é a qualidade percebida do documento de requisitos gerado, avaliada por profissionais da área por meio de uma pesquisa estruturada?
    \item[\textbf{QP4.}] Quão completamente o documento gerado cobre uma especificação de requisitos oficial produzida manualmente (Precisão/Revocação/F1)?
    \item[\textbf{QP5.}] Como os erros se propagam ao longo dos estágios do \textit{pipeline} e como a seleção de modelo por agente afeta o compromisso entre custo e qualidade?
\end{enumerate}

\section{Contribuições}

As principais contribuições deste trabalho são:

\begin{enumerate}
    \item[1.] Um sistema multiagente completo e de código aberto, composto por quatro agentes, para a elicitação de requisitos de ponta a ponta a partir do áudio de reuniões, com suporte a dois formatos de documento (a SRS no padrão IEEE~830 e um formato específico de cliente).
    \item[2.] Uma validação no mundo real sobre um projeto de software de um órgão público, utilizando como referência (\textit{ground truth}) um documento de requisitos oficial produzido manualmente.
    \item[3.] Uma avaliação quantitativa que abrange WER/CER (Agente~1), Precisão/Revocação/F1 com concordância entre avaliadores via \textit{kappa} de Cohen (Agente~2) e uma pesquisa de opinião em escala Likert com profissionais, com avaliação de confiabilidade pelo alfa de Cronbach (Agente~3).
    \item[4.] Uma análise empírica estruturada da propagação de erro entre agentes em um \textit{pipeline} agêntico, acompanhada de medições de latência e de custo de API por estágio.
\end{enumerate}

\section{Organização do Trabalho}

O restante deste trabalho está organizado da seguinte forma. O Capítulo~\ref{cap:fundamentacao} apresenta a fundamentação teórica, revisando os conceitos essenciais ao entendimento do trabalho: engenharia de requisitos e elicitação, histórias de usuário e o \textit{framework} INVEST, a especificação IEEE~830, modelos de linguagem de grande escala, sistemas multiagente e reconhecimento automático de fala. O Capítulo~\ref{cap:relacionados} discute os trabalhos relacionados, posicionando esta pesquisa frente ao estado da arte e destacando as lacunas que a motivam. O Capítulo~\ref{cap:arquitetura} descreve a arquitetura do sistema proposto, o projeto de cada um dos quatro agentes, a estratégia de seleção de modelos e o estudo de caso real. O Capítulo~\ref{cap:avaliacao} apresenta o desenho da avaliação, os resultados obtidos e as ameaças à validade. Por fim, o Capítulo~\ref{cap:conclusao} apresenta as conclusões e os trabalhos futuros.
